\subsection{Models And Detectors}
Each of the five primary models utilizes exactly the same base architecture and the same retraining buffer. Therefore, any variation between them is due to their respective adaptation mechanisms only. The base is an ensemble of a random forest classifier and an extra-trees classifier, both having fifty trees of depth seven. Frozen is trained once on the prefix and receives no subsequent updates. Event-Driven updates when window accuracy falls below a fixed threshold. Full Retraining rebuilds the model at every regime transition. RAPT, RAPT-Enhanced and RAPT-Guard are the controllers of Section~\ref{sec:method}. We also report two baselines that isolate reuse from the trigger: Frozen+Periodic refits on a fixed five-window schedule, and Event-Driven+Reuse uses the event trigger but retrieves a stored policy when the current regime has one. RAPT-Anchor is an oracle upper bound on the gate: each window it predicts with whichever of the active policy and the prefix policy scored better on the labelled lookback, so its inference cost is doubled and charged as such.


ADWIN, Page-Hinkley, EDD and EDMA operate within the same retraining harness, with identical model, buffer, capacity and minimum-retrain-interval. They differ only by how they detect triggers. Following the finding that the historical default parameters left ADWIN and Page-Hinkley silent, we select their operating points on the development split by a fire-rate criterion (median 5--20 events per stream, with the minimum stream event count maximised) rather than against the evaluation streams. The frozen values are ADWIN $\delta=0.2$; Page-Hinkley $\alpha=0.9999$, $\delta=0.005$, threshold $=0.002$, min-instances $=10$; EDD warning $=2$, drift $=3$; EDMA $\alpha=0.3$, $k=1.5$. All four now consume the per-sample 0/1 error stream, which is the input their original definitions assume. Since none of our streams carry ground-truth drift boundary annotations we report detected events, adaptation cost and predictive performance, not detector precision or recall.


\subsection{Metrics And Statistics}
Performance is evaluated through macro-F1, accuracy, precision and recall. Costs are measured two ways. \emph{Adaptation CPU} is the CPU time devoted to adaptation routines (refits, checkpoint writes, and the gate's scoring passes) and \emph{total CPU} is adaptation plus inference, so a method that infers twice is charged twice. A reduction in adaptation time does not translate proportionally into total runtime, which is why we report both; the total-CPU column is what a deployment actually pays. Each experiment uses thirty seeds (42--71) and we report mean and standard deviation over seeds. Comparisons use the Wilcoxon signed-rank test (paired by window, with Holm correction across the comparison family) and Cohen's $d_z$ on the per-seed paired differences, plus two one-sided tests (TOST) for equivalence at a $\pm0.02$ macro-F1 margin. With thirty seeds the seed-level tests are no longer power-limited; the minimum attainable two-sided $p$-value is $2^{-30}$.


\subsection{Predictive Performance}
Table~\ref{tab:primary} holds the primary comparison. On 5G Campus and 5G NR the models cluster within $0.03$ macro-F1, while UGR'16 and NordicDat separate them sharply. On UGR'16 the frozen baseline is the best model, which is the signature of stale-policy reuse: the regimes recur, but the label relation does not, so retrieving the stored policy is worse than never adapting at all.

\begin{table*}[t]
\caption{Primary comparison, mean over 30 seeds. Bold marks the best macro-F1 per stream; CPU in seconds, retrains and reuses are counts. RAPT-Anchor's total CPU includes its doubled inference.}
\label{tab:primary}
\centering
\scriptsize
\begin{tabular}{@{}llrrrrr@{}}
\toprule
\textbf{Dataset} & \textbf{Model} & \textbf{Macro-F1} & \textbf{Adapt CPU} & \textbf{Total CPU} & \textbf{Retr.} & \textbf{Reuse}\\
\midrule
\multirow{7}{*}{Campus} & Frozen & 0.9371 & 0.000 & 1.185 & 0 & 0 \\
 & Event-Driven & 0.9638 & 0.166 & 1.358 & 1 & 0 \\
 & Full Retraining & $\mathbf{0.9836}$ & 2.192 & 3.375 & 14 & 0 \\
 & RAPT & 0.9385 & 0.281 & 1.452 & 2 & 12 \\
 & RAPT-Enhanced & 0.9385 & 0.279 & 1.431 & 2 & 12 \\
 & RAPT-Guard & 0.9391 & 0.499 & 1.675 & 2 & 12 \\
 & RAPT-Anchor & 0.9385 & 0.282 & 3.893 & 2 & 12 \\
\midrule
\multirow{7}{*}{UGR'16} & Frozen & $\mathbf{0.9693}$ & 0.000 & 1.037 & 0 & 0 \\
 & Event-Driven & 0.8978 & 0.835 & 1.942 & 4 & 0 \\
 & Full Retraining & 0.9511 & 29.771 & 30.872 & 144 & 0 \\
 & RAPT & 0.9137 & 2.571 & 3.640 & 11 & 133 \\
 & RAPT-Enhanced & 0.9137 & 2.595 & 3.669 & 11 & 133 \\
 & RAPT-Guard & 0.9357 & 5.743 & 6.810 & 17 & 127 \\
 & RAPT-Anchor & 0.9247 & 2.595 & 5.683 & 11 & 133 \\
\midrule
\multirow{7}{*}{Nordic} & Frozen & 0.2851 & 0.000 & 1.396 & 0 & 0 \\
 & Event-Driven & 0.2739 & 0.346 & 1.706 & 2 & 0 \\
 & Full Retraining & $\mathbf{0.4376}$ & 3.310 & 4.644 & 19 & 0 \\
 & RAPT & 0.3845 & 0.344 & 1.679 & 2 & 16 \\
 & RAPT-Enhanced & 0.3775 & 3.686 & 5.061 & 21 & 16 \\
 & RAPT-Guard & 0.3793 & 1.154 & 2.537 & 5 & 13 \\
 & RAPT-Anchor & 0.3838 & 0.351 & 4.130 & 2 & 16 \\
\midrule
\multirow{7}{*}{5G NR} & Frozen & 0.9110 & 0.000 & 2.709 & 0 & 0 \\
 & Event-Driven & 0.9081 & 1.312 & 3.984 & 11 & 0 \\
 & Full Retraining & 0.9121 & 0.981 & 3.769 & 8 & 0 \\
 & RAPT & 0.8969 & 0.350 & 3.092 & 3 & 5 \\
 & RAPT-Enhanced & 0.9045 & 0.963 & 3.677 & 8 & 5 \\
 & RAPT-Guard & $\mathbf{0.9148}$ & 0.558 & 3.221 & 4 & 4 \\
 & RAPT-Anchor & 0.9042 & 0.349 & 8.851 & 3 & 5 \\
\bottomrule
\end{tabular}
\end{table*}


On 5G Campus every adaptive model is close to saturation. Full Retraining has the best
macro-F1 ($0.9836$), while RAPT reaches $0.9385$; the gap of $0.0451$ is significant
($p = 1.4 \times 10^{-91}$, $d_z = -12.5$). RAPT-Guard edges it up to $0.9391$, still
$0.0445$ below Full Retraining and still significant. Here the gate is not what is
missing: reuse on this stream brings back a policy that is simply less accurate than one
trained on the current buffer, and the gate correctly prefers the prefix model but the
prefix model is also weaker than Full Retraining.

UGR'16 is where the assumption of the mechanism fails, and it teaches us the most. Frozen
has the highest macro-F1 ($0.9693$). Identity-only reuse is the worst of the eight
primary models: RAPT reaches only $0.9137$, which is $0.0556$ below Frozen
($p = 1.0 \times 10^{-68}$, $d_z = -3.6$) and $0.0375$ below Full Retraining, despite its
133 reuses. The twelve regimes recur and their identifiers stay valid across the 30-day
block, but the netflow-feature-to-attack-label mapping does not stay put. RAPT-Guard,
which validates before reusing, recovers to $0.9357$. That is $0.0220$ above base RAPT
and $0.0336$ below Frozen; both gaps are significant ($p = 9.0 \times 10^{-7}$ and
$p = 1.3 \times 10^{-32}$). So the gate repairs most, but not all, of the damage, and on
this stream it does not beat the frozen baseline. RAPT-Enhanced's parity refit reaches
$0.9137$, the same as base RAPT here, because a policy that is poor from the first reused
window is refit late.

NordicDat is hard to predict and strongly imbalanced. Frozen scores $0.2851$, adaptation
helps, and Full Retraining gets $0.4376$ against $0.3845$ for RAPT, a deficit of $0.0531$
($p = 6.1 \times 10^{-4}$, $d_z = -1.1$). RAPT-Guard scores $0.3793$, $0.0583$ below Full
Retraining. Unlike UGR'16, identity-only reuse here is not catastrophic: RAPT stays above
Frozen ($p = 8.4 \times 10^{-73}$), so reuse is safe on this stream even though it is not
accurate. On 5G NR the models are within $0.03$ of one another, and RAPT-Guard has the
best macro-F1 of the RAPT family at $0.9148$. Taken together, the predictive results give
no basis for saying that RAPT improves accuracy. Full Retraining has the best macro-F1 on
NordicDat and 5G NR, the frozen baseline on UGR'16, and a detector baseline on Campus.
RAPT sits below Full Retraining on all four streams. The case for RAPT is cost, and the
case for RAPT-Guard is that it makes reuse safe enough to use where the identity-only
policy is not.

\begin{figure}[t]
\centering
\figorbox{ugr_rolling_accuracy.png}{\columnwidth}{UGR'16 rolling macro-F1 over windows for RAPT, RAPT-Guard and Full Retraining}
\caption{Rolling macro-F1 over evaluation windows on UGR'16. RAPT sags once reused
policies go stale; the prospective validation in RAPT-Guard holds it closer to Full
Retraining.}
\label{fig:ugr_rolling}
\end{figure}


\subsection{Computational Cost and Policy Reuse}
Table~\ref{tab:primary} also lists adaptation and total cost, and Fig.~\ref{fig:cost}
compares total CPU. Relative to Full Retraining the adaptation CPU falls on every stream:
$2.192 \to 0.281$\,s on 5G Campus ($-87.2\%$), $29.771 \to 2.571$\,s on UGR'16
($-91.4\%$), $3.310 \to 0.344$\,s on NordicDat ($-89.6\%$), and $0.981 \to 0.350$\,s on
5G NR ($-64.3\%$). Full Retraining retrains at every regime transition (14, 144, 19 and 8
times); RAPT retrains 2, 11, 2 and 3 times and retrieves a stored policy 12, 133, 16 and
5 times. The gate costs more than identity-only reuse but stays far below Full
Retraining: RAPT-Guard's adaptation CPU is 77.2\%, 80.7\%, 65.1\% and 43.1\% below Full
Retraining on the four streams. The doubled inference of RAPT-Anchor is visible in the
total-CPU column: on 5G NR it totals $8.851$\,s, more than Full Retraining's $3.769$\,s,
even though its adaptation CPU is the lowest of the family. This is why the fair cost
column matters: a method can win on adaptation and lose on total.

\begin{figure}[t]
\centering
\includegraphics[width=\columnwidth]{computational_cost.png}
\caption{Total CPU (adaptation plus inference) per method and stream. RAPT-Anchor's
doubled inference is charged, which lifts it above the other RAPT variants despite its
low adaptation cost.}
\label{fig:cost}
\end{figure}


\subsection{Efficiency Ablation}
\label{sec:ablation}
Table~\ref{tab:ablation} shows the mechanism ladder on 5G
Campus, in which each row adds one mechanism. RAPT-Full is the base controller with the
similarity gate, and it scores $0.9587 \pm 0.0000$ macro-F1 at $0.4283$\,s. Replacing the
absolute refit trigger with a relative-degradation trigger (RAPT-Evidence) leaves
macro-F1 unchanged and fires \emph{zero} refreshes. With an absolute floor (RAPT-Floor),
7.6 refreshes fire and the score is $0.9797 \pm 0.0050$ at $0.5865$\,s, higher than
before but noisier. RAPT-Cheap keeps the periodic schedule and uses the cheaper refit. Its
score is $0.9851 \pm 0.0029$ at $0.8477$\,s, compared with $0.9829 \pm 0.0023$ at
$1.4063$\,s for Full Retraining. The macro-F1 difference is $\Delta = +0.0022$ with
$p = 0.3125$, $d = 0.47$, and because the 95\% interval $[-0.0036, +0.0080]$ includes zero
we cannot claim a difference. RAPT-Incremental scores $0.9549 \pm 0.0020$ at
$0.5813$\,s. Against RAPT-Full it is worse on both axes: lower macro-F1 and higher
adaptation CPU. It is therefore not carried forward. These ablation numbers come from the
legacy 5G-Campus harness and use five seeds, so their seed-level tests are power-limited.

\begin{table}[t]
\caption{Configurations of RAPT on the 5G Campus stream, five seeds.}
\label{tab:ablation}
\centering
\footnotesize
\begin{tabular}{@{}lccc@{}}
\toprule
{\bf Configuration} & {\bf Macro-F1} & {\bf Adapt.\ CPU} & {\bf Refreshes}\\
\midrule
Full Retrain. (ref.) & $0.9829 \pm 0.0023$& $1.4063$ & -- \\
RAPT-Full & $0.9587 \pm 0.0000$ & $0.4283$ & 0 \\
RAPT-Evidence & $0.9587 \pm 0.0000$ & $0.4260$ & 0\\
RAPT-Floor & $0.9797 \pm 0.0050$ & $0.5865$ & 7.6 \\
RAPT-Cheap & $\mathbf{0.9851 \pm 0.0029}$ & $0.8477$ & 22 \\
RAPT-Incremental & $0.9549 \pm 0.0020$ & $0.5813$ & 0\\
\bottomrule
\end{tabular}
\end{table}


\subsection{Determinants of Cost and Trigger Behavior}
\label{sec:cost}
The cheap refresh was chosen from direct refit-timing measurements, varying one factor at a time with the ensemble held fixed: a 20-tree refit takes about 22\,ms against 104\,ms for 100 trees, and a 200-observation buffer about 80\,ms against 104\,ms for 1000 observations. Tree count dominates, so the cheap refit cuts refit time by about four-fifths at 20 trees. The zero-refresh result for RAPT-Evidence is a property of the mechanism, not a tuning outcome: a relative trigger compares accuracy against a decayed reference built from that policy's own earlier accuracy, so when a regime is persistently poor the reference decays toward the poor level and the gap the rule tests for never opens. RAPT-Guard sidesteps this by testing the stored policy directly against the prefix policy at the moment of reuse.


\subsection{Validation-Gate Accuracy}
\label{sec:guard}
RAPT-Guard turns every regime return into a decision, so we can audit that decision
directly. For each return we compare the gate's choice with the outcome on the incoming
window, using a two-point tolerance: a reuse is counted as harmful only if the stored
policy scores more than two points below the prefix policy on that window, and a block is
counted as a miss only if reuse would have been more than two points better. Table~\ref{tab:gate}
reports the counts. The gate almost never allows a harmful reuse: 0 of 36 on Campus, 6 of
385 on UGR'16 ($1.6\%$), 3 of 39 on NordicDat ($7.7\%$), and 0 of 11 on 5G NR. It is
slightly conservative on NordicDat, blocking 3 reuses that would have helped. The margin
itself was chosen on the development split, and the frozen margin is the one that keeps
the gate a gate (allow rate above $0.30$) while maximising development macro-F1: at a
positive margin the gate blocks almost everything, which is not a gate, and at a negative
margin it allows everything, which is not safe. The zero margin is the corner that still
reuses on three of the four streams and blocks most harmful reuses on the fourth.

\begin{table}[t]
\caption{RAPT-Guard decisions at every regime return, three seeds, with a two-point tolerance. Bad allow counts a reuse that lost more than two macro-F1 points to the prefix policy on the incoming window.}
\label{tab:gate}
\centering
\footnotesize
\begin{tabular}{@{}lrrrrrr@{}}
\toprule
\textbf{Stream} & \textbf{Returns} & \textbf{Allow} & \textbf{Block} & \textbf{Bad allow} & \textbf{Bad block} & \textbf{Bad-allow rate}\\
\midrule
Campus & 36 & 36 & 0 & 0 & 0 & 0.000 \\
UGR'16 & 399 & 385 & 14 & 6 & 0 & 0.016 \\
Nordic & 48 & 39 & 9 & 3 & 3 & 0.077 \\
5G NR & 15 & 11 & 4 & 0 & 0 & 0.000 \\
\bottomrule
\end{tabular}
\end{table}



\subsection{Comparison to Historical Detectors}
\label{sec:detectors}
Table~\ref{tab:detectors} and Fig.~\ref{fig:detector} report the four historical
detectors at the frozen development operating points. They no longer reduce to the frozen
baseline. Page-Hinkley is the strongest detector on Campus and NordicDat ($0.9843$ and
$0.4676$), and matches EDD and EDMA on UGR'16 ($0.9598$), but it fires $6$--$34$ times per
stream and pays the most adaptation CPU of any method ($2.47$--$5.44$\,s). ADWIN fires
$8.7$ times on UGR'16 and $30$ times on NordicDat yet has the weakest NordicDat macro-F1
($0.328$), so firing more is not the same as adapting well. RAPT fires $2$--$11$ times and
RAPT-Guard $4$--$17$ times, both well below the detector baselines, at a fraction of their
adaptation cost. On Campus the detector baselines reach Full-Retraining accuracy at lower
cost, so RAPT is not the accuracy leader there; on UGR'16 the frozen baseline still wins.
RAPT remains the lowest-cost adaptive option that stays near the frozen baseline where
reuse is safe, and RAPT-Guard is the option that keeps reuse safe where it is not.

\begin{table}[t]
\caption{Historical detectors and the reuse methods at the frozen operating points, mean over 30 seeds. Each cell is macro-F1 with retrain events per seed in parentheses.}
\label{tab:detectors}
\centering
\footnotesize
\begin{tabular}{@{}lrrrr@{}}
\toprule
\textbf{Method} & \textbf{Campus} & \textbf{UGR'16} & \textbf{Nordic} & \textbf{5G NR}\\
\midrule
ADWIN & 0.937 (0.0) & 0.950 (8.7) & 0.328 (30.5) & 0.911 (0.1) \\
Page-Hinkley & 0.984 (6.0) & 0.960 (13.3) & 0.468 (34.1) & 0.905 (19.6) \\
EDD & 0.984 (6.0) & 0.960 (13.3) & 0.433 (32.1) & 0.910 (7.4) \\
EDMA & 0.984 (6.0) & 0.960 (13.3) & 0.378 (29.6) & 0.907 (9.0) \\
Event-Driven & 0.964 (1.0) & 0.898 (3.6) & 0.274 (2.0) & 0.908 (10.9) \\
RAPT & 0.939 (2.0) & 0.914 (11.0) & 0.385 (2.0) & 0.897 (3.0) \\
\bottomrule
\end{tabular}
\end{table}



\begin{figure}[t]
\centering
\includegraphics[width=\columnwidth]{detector_comparison.png}
\caption{Retrain events per seed for the historical detectors and the reuse methods,
mean over 30 seeds. ADWIN and Page-Hinkley now fire at sensible rates; RAPT and
RAPT-Guard fire least.}
\label{fig:detector}
\end{figure}


\subsection{Gate Cadence: Short versus Long Window}
\label{sec:cadence}
The gate can be evaluated at two cadences: on the incoming window only (a short,
boundary-only lookback) or on every window while a reused policy is active (a long,
continuous lookback). We test both under a synthetic concept-drift stress panel that
reverses the rank order of the top predictive features inside a regime at three
severities. Under the boundary-only cadence no refit ever fires, because the drift is
in-regime and there is no boundary to key on; under the every-window cadence the gate
fires $7$--$10$ refits at the two higher severities. The macro-F1 effect is small at the
severities where the underlying window is still predictable ($0.8902 \to 0.8941$ at
severity $0.5$, no change at $0.2$ and $1.0$), so the honest reading is that the
long-window cadence detects in-regime drift that the short-window cadence cannot see at
all, and recovers a little accuracy on the mid-severity panel, but neither cadence moves
macro-F1 much once the window itself becomes near-single-class. We keep the long-window
gate as a detection mechanism and do not claim a large accuracy effect.
